% =============================================================
%  MAT231-F27 Calculus II — Lecture 18
%  Unit 18: Double integrals over rectangles and general regions;
%  Fubini's theorem
%
%  Build from the repository root:   npm run materials -- lecture-18
%  This writes lecture-18.pdf (slides, with overlays) and
%  lecture-18-handout.pdf (print version) into this folder.
% =============================================================
\documentclass[10pt]{wosm-slides}

\usepackage{pgfplots}
\pgfplotsset{compat=1.18}
\pgfplotsset{wosm axis/.style={
  axis lines=left, axis line style={wosm-muted, -},
  tick style={wosm-muted}, tick label style={font=\scriptsize, text=wosm-muted},
  label style={font=\scriptsize, text=wosm-muted},
  legend style={draw=none, fill=none, font=\scriptsize, text=wosm-text},
  legend cell align=left}}

\newcommand{\R}{\mathbb{R}}
\newcommand{\dA}{\, dA}

\course{MAT231-F27}{Calculus II}
\title[L18 \textperiodcentered\ Double integrals]{Double integrals}
\subtitle{Rectangles, general regions, and Fubini's theorem}
\author{Dr.\ Ali Khudiyev}
\institute{WoSM Academy}
\date{Lecture 18\enspace\textperiodcentered\enspace Unit 18\enspace\textperiodcentered\enspace Fall 2027}

\begin{document}

\maketitle

% -------------------------------------------------------------
\begin{frame}{Last time}
  \begin{itemize}
    \item Implicit function theorem: $F(\mathbf x, \mathbf y) = \mathbf c$ can be solved locally for
      $\mathbf y$ when $D_{\mathbf y}F$ is invertible.
    \item Inverse function theorem: $\det DF(\mathbf a) \ne 0$ gives a local inverse.
    \item This completes differential calculus in several variables.
  \end{itemize}
  \vskip0.6em
  Now integration. In one variable, $\int_a^b f(x)\,dx$ is the area under a graph. Today: the volume
  under a surface $z = f(x, y)$, and how to compute it with one-variable integrals.
\end{frame}

\begin{frame}{This lecture}
  \tableofcontents
  \vskip1em
  {\small\color{wosm-muted}
  Reading: Strang \& Herman, \emph{Calculus, Volume 3}, Sections 5.1 and 5.2.\par}
\end{frame}

% =============================================================
\section{Double integrals over rectangles}

\begin{frame}{Volume under a surface}
  \begin{columns}[c]
    \begin{column}{0.52\textwidth}
      Let $f \ge 0$ on the rectangle $R = [a, b] \times [c, d]$. To find the volume under its graph:
      \begin{enumerate}
        \item cut $R$ into $m \times n$ small rectangles $R_{ij}$, each of area
          $\Delta A = \Delta x\, \Delta y$;
        \item pick a sample point $(x_{ij}^*, y_{ij}^*)$ in each;
        \item add up the volumes of the boxes:
          \[ \sum_{i=1}^m \sum_{j=1}^n f(x_{ij}^*, y_{ij}^*)\, \Delta A . \]
      \end{enumerate}
      This is a \textbf{Riemann sum}.
    \end{column}
    \begin{column}{0.46\textwidth}
      \centering
      \begin{tikzpicture}[x={(-0.62cm,-0.36cm)}, y={(1.15cm,-0.1cm)}, z={(0cm,1cm)}]
        \draw[wosm-muted, -latex] (0,0,0) -- (2.6,0,0) node[below left, font=\scriptsize] {$x$};
        \draw[wosm-muted, -latex] (0,0,0) -- (0,2.7,0) node[right, font=\scriptsize] {$y$};
        \draw[wosm-muted, -latex] (0,0,0) -- (0,0,2.9) node[above, font=\scriptsize] {$z$};
        \draw[wosm-accent!45!wosm-bg] plot coordinates {(0.000,0.000,2.400) (0.000,0.100,2.397) (0.000,0.200,2.388) (0.000,0.300,2.373) (0.000,0.400,2.352) (0.000,0.500,2.325) (0.000,0.600,2.292) (0.000,0.700,2.253) (0.000,0.800,2.208) (0.000,0.900,2.157) (0.000,1.000,2.100) (0.000,1.100,2.037) (0.000,1.200,1.968) (0.000,1.300,1.893) (0.000,1.400,1.812) (0.000,1.500,1.725) (0.000,1.600,1.632) (0.000,1.700,1.533) (0.000,1.800,1.428) (0.000,1.900,1.317) (0.000,2.000,1.200)};
        \draw[wosm-accent!45!wosm-bg] plot coordinates {(0.000,0.000,2.400) (0.100,0.000,2.397) (0.200,0.000,2.388) (0.300,0.000,2.373) (0.400,0.000,2.352) (0.500,0.000,2.325) (0.600,0.000,2.292) (0.700,0.000,2.253) (0.800,0.000,2.208) (0.900,0.000,2.157) (1.000,0.000,2.100) (1.100,0.000,2.037) (1.200,0.000,1.968) (1.300,0.000,1.893) (1.400,0.000,1.812) (1.500,0.000,1.725) (1.600,0.000,1.632) (1.700,0.000,1.533) (1.800,0.000,1.428) (1.900,0.000,1.317) (2.000,0.000,1.200)};
        \draw[wosm-accent!45!wosm-bg] plot coordinates {(0.500,0.000,2.325) (0.500,0.100,2.322) (0.500,0.200,2.313) (0.500,0.300,2.298) (0.500,0.400,2.277) (0.500,0.500,2.250) (0.500,0.600,2.217) (0.500,0.700,2.178) (0.500,0.800,2.133) (0.500,0.900,2.082) (0.500,1.000,2.025) (0.500,1.100,1.962) (0.500,1.200,1.893) (0.500,1.300,1.818) (0.500,1.400,1.737) (0.500,1.500,1.650) (0.500,1.600,1.557) (0.500,1.700,1.458) (0.500,1.800,1.353) (0.500,1.900,1.242) (0.500,2.000,1.125)};
        \draw[wosm-accent!45!wosm-bg] plot coordinates {(0.000,0.500,2.325) (0.100,0.500,2.322) (0.200,0.500,2.313) (0.300,0.500,2.298) (0.400,0.500,2.277) (0.500,0.500,2.250) (0.600,0.500,2.217) (0.700,0.500,2.178) (0.800,0.500,2.133) (0.900,0.500,2.082) (1.000,0.500,2.025) (1.100,0.500,1.962) (1.200,0.500,1.893) (1.300,0.500,1.818) (1.400,0.500,1.737) (1.500,0.500,1.650) (1.600,0.500,1.557) (1.700,0.500,1.458) (1.800,0.500,1.353) (1.900,0.500,1.242) (2.000,0.500,1.125)};
        \draw[wosm-accent!45!wosm-bg] plot coordinates {(1.000,0.000,2.100) (1.000,0.100,2.097) (1.000,0.200,2.088) (1.000,0.300,2.073) (1.000,0.400,2.052) (1.000,0.500,2.025) (1.000,0.600,1.992) (1.000,0.700,1.953) (1.000,0.800,1.908) (1.000,0.900,1.857) (1.000,1.000,1.800) (1.000,1.100,1.737) (1.000,1.200,1.668) (1.000,1.300,1.593) (1.000,1.400,1.512) (1.000,1.500,1.425) (1.000,1.600,1.332) (1.000,1.700,1.233) (1.000,1.800,1.128) (1.000,1.900,1.017) (1.000,2.000,0.900)};
        \draw[wosm-accent!45!wosm-bg] plot coordinates {(0.000,1.000,2.100) (0.100,1.000,2.097) (0.200,1.000,2.088) (0.300,1.000,2.073) (0.400,1.000,2.052) (0.500,1.000,2.025) (0.600,1.000,1.992) (0.700,1.000,1.953) (0.800,1.000,1.908) (0.900,1.000,1.857) (1.000,1.000,1.800) (1.100,1.000,1.737) (1.200,1.000,1.668) (1.300,1.000,1.593) (1.400,1.000,1.512) (1.500,1.000,1.425) (1.600,1.000,1.332) (1.700,1.000,1.233) (1.800,1.000,1.128) (1.900,1.000,1.017) (2.000,1.000,0.900)};
        \draw[wosm-accent!45!wosm-bg] plot coordinates {(1.500,0.000,1.725) (1.500,0.100,1.722) (1.500,0.200,1.713) (1.500,0.300,1.698) (1.500,0.400,1.677) (1.500,0.500,1.650) (1.500,0.600,1.617) (1.500,0.700,1.578) (1.500,0.800,1.533) (1.500,0.900,1.482) (1.500,1.000,1.425) (1.500,1.100,1.362) (1.500,1.200,1.293) (1.500,1.300,1.218) (1.500,1.400,1.137) (1.500,1.500,1.050) (1.500,1.600,0.957) (1.500,1.700,0.858) (1.500,1.800,0.753) (1.500,1.900,0.642) (1.500,2.000,0.525)};
        \draw[wosm-accent!45!wosm-bg] plot coordinates {(0.000,1.500,1.725) (0.100,1.500,1.722) (0.200,1.500,1.713) (0.300,1.500,1.698) (0.400,1.500,1.677) (0.500,1.500,1.650) (0.600,1.500,1.617) (0.700,1.500,1.578) (0.800,1.500,1.533) (0.900,1.500,1.482) (1.000,1.500,1.425) (1.100,1.500,1.362) (1.200,1.500,1.293) (1.300,1.500,1.218) (1.400,1.500,1.137) (1.500,1.500,1.050) (1.600,1.500,0.957) (1.700,1.500,0.858) (1.800,1.500,0.753) (1.900,1.500,0.642) (2.000,1.500,0.525)};
        \draw[wosm-accent!45!wosm-bg] plot coordinates {(2.000,0.000,1.200) (2.000,0.100,1.197) (2.000,0.200,1.188) (2.000,0.300,1.173) (2.000,0.400,1.152) (2.000,0.500,1.125) (2.000,0.600,1.092) (2.000,0.700,1.053) (2.000,0.800,1.008) (2.000,0.900,0.957) (2.000,1.000,0.900) (2.000,1.100,0.837) (2.000,1.200,0.768) (2.000,1.300,0.693) (2.000,1.400,0.612) (2.000,1.500,0.525) (2.000,1.600,0.432) (2.000,1.700,0.333) (2.000,1.800,0.228) (2.000,1.900,0.117) (2.000,2.000,0.000)};
        \draw[wosm-accent!45!wosm-bg] plot coordinates {(0.000,2.000,1.200) (0.100,2.000,1.197) (0.200,2.000,1.188) (0.300,2.000,1.173) (0.400,2.000,1.152) (0.500,2.000,1.125) (0.600,2.000,1.092) (0.700,2.000,1.053) (0.800,2.000,1.008) (0.900,2.000,0.957) (1.000,2.000,0.900) (1.100,2.000,0.837) (1.200,2.000,0.768) (1.300,2.000,0.693) (1.400,2.000,0.612) (1.500,2.000,0.525) (1.600,2.000,0.432) (1.700,2.000,0.333) (1.800,2.000,0.228) (1.900,2.000,0.117) (2.000,2.000,0.000)};
        \filldraw[fill=wosm-surface, draw=wosm-text, thin] (0.000,1.000,0.000) -- (1.000,1.000,0.000) -- (1.000,1.000,2.250) -- (0.000,1.000,2.250) -- cycle;
        \filldraw[fill=wosm-rule, draw=wosm-text, thin] (1.000,0.000,0.000) -- (1.000,1.000,0.000) -- (1.000,1.000,2.250) -- (1.000,0.000,2.250) -- cycle;
        \filldraw[fill=wosm-accent!25!wosm-bg, draw=wosm-text, thin] (0.000,0.000,2.250) -- (1.000,0.000,2.250) -- (1.000,1.000,2.250) -- (0.000,1.000,2.250) -- cycle;
        \fill[wosm-brick] (0.500,0.500,2.250) circle (1.2pt);
        \filldraw[fill=wosm-surface, draw=wosm-text, thin] (0.000,2.000,0.000) -- (1.000,2.000,0.000) -- (1.000,2.000,1.650) -- (0.000,2.000,1.650) -- cycle;
        \filldraw[fill=wosm-rule, draw=wosm-text, thin] (1.000,1.000,0.000) -- (1.000,2.000,0.000) -- (1.000,2.000,1.650) -- (1.000,1.000,1.650) -- cycle;
        \filldraw[fill=wosm-accent!25!wosm-bg, draw=wosm-text, thin] (0.000,1.000,1.650) -- (1.000,1.000,1.650) -- (1.000,2.000,1.650) -- (0.000,2.000,1.650) -- cycle;
        \fill[wosm-brick] (0.500,1.500,1.650) circle (1.2pt);
        \filldraw[fill=wosm-surface, draw=wosm-text, thin] (1.000,1.000,0.000) -- (2.000,1.000,0.000) -- (2.000,1.000,1.650) -- (1.000,1.000,1.650) -- cycle;
        \filldraw[fill=wosm-rule, draw=wosm-text, thin] (2.000,0.000,0.000) -- (2.000,1.000,0.000) -- (2.000,1.000,1.650) -- (2.000,0.000,1.650) -- cycle;
        \filldraw[fill=wosm-accent!25!wosm-bg, draw=wosm-text, thin] (1.000,0.000,1.650) -- (2.000,0.000,1.650) -- (2.000,1.000,1.650) -- (1.000,1.000,1.650) -- cycle;
        \fill[wosm-brick] (1.500,0.500,1.650) circle (1.2pt);
        \filldraw[fill=wosm-surface, draw=wosm-text, thin] (1.000,2.000,0.000) -- (2.000,2.000,0.000) -- (2.000,2.000,1.050) -- (1.000,2.000,1.050) -- cycle;
        \filldraw[fill=wosm-rule, draw=wosm-text, thin] (2.000,1.000,0.000) -- (2.000,2.000,0.000) -- (2.000,2.000,1.050) -- (2.000,1.000,1.050) -- cycle;
        \filldraw[fill=wosm-accent!25!wosm-bg, draw=wosm-text, thin] (1.000,1.000,1.050) -- (2.000,1.000,1.050) -- (2.000,2.000,1.050) -- (1.000,2.000,1.050) -- cycle;
        \fill[wosm-brick] (1.500,1.500,1.050) circle (1.2pt);
      \end{tikzpicture}
      \par\small $f = 8 - x^2 - y^2$ on $[0, 2]^2$, with $2 \times 2$ boxes at the midpoints.
    \end{column}
  \end{columns}
\end{frame}

\begin{frame}{The double integral}
  \begin{definition}
    $\displaystyle \iint_R f(x, y) \dA = \lim_{m, n \to \infty} \sum_{i=1}^m \sum_{j=1}^n f(x_{ij}^*, y_{ij}^*)\, \Delta A$, \
    if the limit exists and is the same for every choice of sample points. Then $f$ is
    \textbf{integrable} on $R$.
  \end{definition}
  \pause
  \begin{theorem}
    A continuous function on $R$ is integrable. So is a bounded function that is continuous except on
    finitely many smooth curves. (Proofs as in Calculus~I, Units 18--19.)
  \end{theorem}
  \pause
  For $f \ge 0$: the volume under the graph. In general: the volume above the $xy$-plane minus the
  volume below.
\end{frame}

\begin{frame}{Example: Riemann sums}
  $f(x, y) = 8 - x^2 - y^2$ on $R = [0, 2] \times [0, 2]$.
  \pause
  \begin{itemize}
    \item $2 \times 2$ squares ($\Delta A = 1$), upper right corners $(1, 1), (1, 2), (2, 1), (2, 2)$:
      \[ (6 + 3 + 3 + 0) \cdot 1 = 12 . \]
    \pause
    \item $2 \times 2$ squares, midpoints $(\frac12, \frac12), (\frac12, \frac32), (\frac32, \frac12), (\frac32, \frac32)$:
      \[ (7.5 + 5.5 + 5.5 + 3.5) \cdot 1 = 22 . \]
    \pause
    \item $4 \times 4$ squares ($\Delta A = \frac14$), midpoints: \ $21.5$.
  \end{itemize}
  \pause
  The exact value is $\frac{64}{3} \approx 21.33$ (shortly). The corners underestimate because $f$
  decreases in $x$ and in $y$; the midpoints do much better, as in one variable.
\end{frame}

\begin{frame}{Properties}
  For integrable $f$, $g$ on $R$ and constants $\alpha$, $\beta$:
  \begin{itemize}
    \item \textbf{Linearity}: $\iint_R (\alpha f + \beta g) \dA = \alpha \iint_R f \dA + \beta \iint_R g \dA$.
    \item \textbf{Monotonicity}: if $f \le g$ on $R$, then $\iint_R f \dA \le \iint_R g \dA$.
    \item \textbf{Area}: $\iint_R 1 \dA = \operatorname{area}(R)$. \ So if $m \le f \le M$, then
      \[ m \cdot \operatorname{area}(R) \le \iint_R f \dA \le M \cdot \operatorname{area}(R) . \]
    \item \textbf{Additivity}: if $R$ is cut into two rectangles $R_1$, $R_2$, then
      $\iint_R f = \iint_{R_1} f + \iint_{R_2} f$.
  \end{itemize}
  \pause
  The \textbf{average value} of $f$ on $R$ is $\displaystyle \bar f = \frac{1}{\operatorname{area}(R)} \iint_R f \dA$.
\end{frame}

% =============================================================
\section{Iterated integrals and Fubini's theorem}

\begin{frame}{Computing by slicing}
  Fix $x$, and integrate over $y$:
  \[ A(x) = \int_c^d f(x, y)\, dy \]
  is the area of the slice of the solid at $x$. Adding up the slices (as for volumes in
  Calculus~I, Unit 24):
  \[ \text{volume} = \int_a^b A(x)\, dx = \int_a^b \Bigl( \int_c^d f(x, y)\, dy \Bigr) dx . \]
  \pause
  This is an \textbf{iterated integral}, usually written $\int_a^b \int_c^d f(x, y)\, dy\, dx$: the
  inner integral belongs to the inner differential. Slicing the other way gives
  $\int_c^d \int_a^b f(x, y)\, dx\, dy$.
\end{frame}

\begin{frame}{Fubini's theorem}
  \begin{theorem}[Fubini]
    If $f$ is continuous on $R = [a, b] \times [c, d]$, then
    \[ \iint_R f \dA = \int_a^b \int_c^d f(x, y)\, dy\, dx = \int_c^d \int_a^b f(x, y)\, dx\, dy . \]
  \end{theorem}
  \pause
  \textbf{Proof idea.} Group a Riemann sum by columns:
  \[ \sum_i \Bigl( \sum_j f(x_i^*, y_j^*)\, \Delta y \Bigr) \Delta x . \]
  The inner sum approximates $A(x_i^*) = \int_c^d f(x_i^*, y)\, dy$, and the outer one is then a Riemann
  sum for $\int_a^b A(x)\, dx$. Grouping by rows gives the other order.
\end{frame}

\begin{frame}{Examples}
  \begin{example}[The volume from before]
    $\displaystyle \int_0^2 \int_0^2 (8 - x^2 - y^2)\, dy\, dx = \int_0^2 \Bigl(16 - 2x^2 - \frac83\Bigr) dx
      = 32 - \frac{16}{3} - \frac{16}{3} = \frac{64}{3}$.
  \end{example}
  \pause
  \begin{example}[A product]
    If $f(x, y) = g(x)\, h(y)$, the integral splits:
    $\iint_R g(x) h(y) \dA = \Bigl(\int_a^b g\Bigr)\Bigl(\int_c^d h\Bigr)$. For instance
    \[ \int_0^1 \int_0^\pi x^2 \sin y\, dy\, dx = \Bigl(\int_0^1 x^2\, dx\Bigr) \Bigl(\int_0^\pi \sin y\, dy\Bigr)
       = \frac13 \cdot 2 = \frac23 . \]
  \end{example}
\end{frame}

\begin{frame}{Choose the order wisely}
  \begin{example}[{$\iint_R x\, e^{xy} \dA$, \ $R = [0, 1] \times [0, 1]$}]
    \textbf{$y$ first}: for fixed $x$, an antiderivative of $x e^{xy}$ in $y$ is $e^{xy}$:
    \[ \int_0^1 \int_0^1 x e^{xy}\, dy\, dx = \int_0^1 (e^x - 1)\, dx = e - 2 . \]
    \pause
    \textbf{$x$ first}: $\int_0^1 x e^{xy}\, dx$ needs integration by parts, and then the $y$-integral is
    harder. Same answer, more work.
  \end{example}
  \pause
  {\small Fubini guarantees the two orders agree; it is our job to pick the easier one.\par}
\end{frame}

\begin{frame}{When Fubini fails}
  \begin{example}[{$f(x, y) = \dfrac{x^2 - y^2}{(x^2 + y^2)^2}$ on $[0, 1] \times [0, 1]$}]
    Since $\frac{\partial}{\partial y}\, \frac{y}{x^2 + y^2} = f(x, y)$, \
    $\displaystyle \int_0^1 \int_0^1 f\, dy\, dx = \int_0^1 \frac{dx}{1 + x^2} = \frac\pi4$. \
    By the symmetry $f(y, x) = -f(x, y)$, \
    $\displaystyle \int_0^1 \int_0^1 f\, dx\, dy = -\frac\pi4$.
  \end{example}
  \pause
  \begin{warning}[Check the hypothesis]
    $f$ is not continuous on the closed square: it is unbounded near $(0, 0)$, and not integrable.
    Then the two iterated integrals need not agree.
  \end{warning}
\end{frame}

% =============================================================
\section{General regions}

\begin{frame}{Integrating over a general region}
  Let $D$ be a bounded region, inside a rectangle $R$. Extend $f$ by zero:
  \[ F(x, y) = \begin{cases} f(x, y), & (x, y) \in D, \\ 0, & (x, y) \in R \setminus D, \end{cases}
     \qquad \iint_D f \dA = \iint_R F \dA . \]
  \pause
  \begin{itemize}
    \item If $f$ is continuous on $D$ and the boundary of $D$ consists of finitely many smooth curves,
      $F$ is integrable (its discontinuities lie on those curves).
    \item $\iint_D 1 \dA = \operatorname{area}(D)$.
    \item Additivity: if $D = D_1 \cup D_2$ and $D_1$, $D_2$ overlap only along their boundaries,
      $\iint_D f = \iint_{D_1} f + \iint_{D_2} f$.
  \end{itemize}
\end{frame}

\begin{frame}{Type I and type II regions}
  \begin{columns}[t]
    \begin{column}{0.49\textwidth}
      \textbf{Type I}: \ $a \le x \le b$, \ $g_1(x) \le y \le g_2(x)$.
      \[ \iint_D f \dA = \int_a^b \int_{g_1(x)}^{g_2(x)} f(x, y)\, dy\, dx \]
      \centering
      \begin{tikzpicture}[scale=2.2]
        \fill[wosm-accent!18!wosm-bg] plot[domain=0:1, samples=40] (\x, \x*\x) -- (1,1) -- (0,0);
        \fill[wosm-accent!55!wosm-bg] (0.55,0.33) rectangle (0.62,0.585);
        \draw[wosm-text, thick, domain=0:1, samples=40] plot (\x, \x*\x);
        \draw[wosm-text, thick] (0,0) -- (1,1);
        \node[font=\scriptsize, below right] at (0.78,0.6) {$y = x^2$};
        \node[font=\scriptsize, above left] at (0.62,0.62) {$y = x$};
        \draw[wosm-muted, -latex] (-0.1,0) -- (1.2,0);
        \draw[wosm-muted, -latex] (0,-0.1) -- (0,1.15);
      \end{tikzpicture}
      \par\small Vertical strips: $y$ runs from the lower curve to the upper one.
    \end{column}
    \begin{column}{0.49\textwidth}
      \textbf{Type II}: \ $c \le y \le d$, \ $h_1(y) \le x \le h_2(y)$.
      \[ \iint_D f \dA = \int_c^d \int_{h_1(y)}^{h_2(y)} f(x, y)\, dx\, dy \]
      \centering
      \begin{tikzpicture}[scale=2.2]
        \fill[wosm-accent!18!wosm-bg] plot[domain=0:1, samples=40] (\x, \x*\x) -- (1,1) -- (0,0);
        \fill[wosm-accent!55!wosm-bg] (0.45,0.45) rectangle (0.6708,0.52);
        \draw[wosm-text, thick, domain=0:1, samples=40] plot (\x, \x*\x);
        \draw[wosm-text, thick] (0,0) -- (1,1);
        \node[font=\scriptsize, below right] at (0.78,0.6) {$x = \sqrt y$};
        \node[font=\scriptsize, above left] at (0.4,0.4) {$x = y$};
        \draw[wosm-muted, -latex] (-0.1,0) -- (1.2,0);
        \draw[wosm-muted, -latex] (0,-0.1) -- (0,1.15);
      \end{tikzpicture}
      \par\small Horizontal strips: $x$ runs from the left curve to the right one.
    \end{column}
  \end{columns}
\end{frame}

\begin{frame}{Example: both ways}
  $D$ is the region between $y = x^2$ and $y = x$ (the picture), and $f(x, y) = xy$.
  \pause
  \begin{itemize}
    \item \textbf{Type I}: $0 \le x \le 1$, $x^2 \le y \le x$:
      \[ \int_0^1 \int_{x^2}^{x} xy\, dy\, dx = \int_0^1 x \cdot \frac{x^2 - x^4}{2}\, dx
         = \frac12 \Bigl(\frac14 - \frac16\Bigr) = \frac{1}{24} . \]
    \pause
    \item \textbf{Type II}: $0 \le y \le 1$, $y \le x \le \sqrt y$:
      \[ \int_0^1 \int_{y}^{\sqrt y} xy\, dx\, dy = \int_0^1 y \cdot \frac{y - y^2}{2}\, dy
         = \frac12 \Bigl(\frac13 - \frac14\Bigr) = \frac{1}{24} . \]
  \end{itemize}
  \pause
  {\small And $\operatorname{area}(D) = \int_0^1 \int_{x^2}^x 1\, dy\, dx = \int_0^1 (x - x^2)\, dx = \frac16$.\par}
\end{frame}

\begin{frame}{Example: a volume}
  \begin{example}[The tetrahedron under $z = 1 - x - y$, above the triangle $x, y \ge 0$, $x + y \le 1$]
    The triangle is of type I with $0 \le x \le 1$, $0 \le y \le 1 - x$:
    \[ V = \int_0^1 \int_0^{1-x} (1 - x - y)\, dy\, dx
         = \int_0^1 \frac{(1 - x)^2}{2}\, dx = \frac16 . \]
  \end{example}
  \pause
  {\small Check: a pyramid has volume $\frac13 \cdot \text{base} \cdot \text{height} = \frac13 \cdot \frac12 \cdot 1 = \frac16$.\par}
\end{frame}

\begin{frame}{Setting up the limits}
  \begin{block}{For a type I description}
    \begin{enumerate}
      \item Sketch $D$. The outer limits are the smallest and largest $x$ in $D$.
      \item For fixed $x$, the vertical line enters and leaves $D$ at the inner limits for $y$.
      \item If the entry or exit curve changes formula, split $D$ into pieces.
    \end{enumerate}
  \end{block}
  \pause
  \begin{example}[The triangle with vertices $(0, 0)$, $(2, 0)$, $(1, 1)$]
    Type I needs two pieces (the top edge changes at $x = 1$); type II needs one:
    \[ \int_0^1 \!\! \int_0^x f\, dy\, dx + \int_1^2 \!\! \int_0^{2-x} f\, dy\, dx
       \ = \ \int_0^1 \!\! \int_y^{2-y} f\, dx\, dy . \]
  \end{example}
\end{frame}

\begin{frame}{Changing the order of integration}
  \begin{columns}[c]
    \begin{column}{0.62\textwidth}
      $\displaystyle \int_0^1 \int_x^1 e^{y^2}\, dy\, dx$ \ --- \ $e^{y^2}$ has no elementary antiderivative.
      \pause
      \vskip0.5em
      The region: $0 \le x \le 1$, $x \le y \le 1$, the triangle $0 \le x \le y \le 1$. As a type II
      region: $0 \le y \le 1$, $0 \le x \le y$. So
      \[ \int_0^1 \int_0^y e^{y^2}\, dx\, dy = \int_0^1 y\, e^{y^2}\, dy = \frac{e - 1}{2} . \]
      \pause
      {\small The inner integral over $x$ produced the factor $y$ that the substitution
      $u = y^2$ needs.\par}
    \end{column}
    \begin{column}{0.36\textwidth}
      \centering
      \begin{tikzpicture}[scale=2.6]
        \fill[wosm-accent!18!wosm-bg] (0,0) -- (1,1) -- (0,1) -- cycle;
        \fill[wosm-accent!55!wosm-bg] (0.3,0.3) rectangle (0.36,1);
        \fill[wosm-brick!35!wosm-bg] (0,0.7) rectangle (0.7,0.75);
        \draw[wosm-text, thick] (0,0) -- (1,1) node[right, font=\scriptsize] {$y = x$};
        \draw[wosm-text, thick] (0,1) -- (1,1);
        \draw[wosm-muted, -latex] (-0.1,0) -- (1.15,0) node[right, font=\scriptsize] {$x$};
        \draw[wosm-muted, -latex] (0,-0.1) -- (0,1.15) node[above, font=\scriptsize] {$y$};
      \end{tikzpicture}
      \par\small Vertical strip: $y$ from $x$ to $1$. \\ Horizontal strip: $x$ from $0$ to $y$.
    \end{column}
  \end{columns}
\end{frame}

% =============================================================
\section{Applications}

\begin{frame}{Mass and centre of mass}
  A thin plate (\textbf{lamina}) occupying $D$, with density $\rho(x, y)$ (mass per unit area):
  \[ M = \iint_D \rho \dA, \qquad
     \bar x = \frac1M \iint_D x\, \rho \dA, \qquad
     \bar y = \frac1M \iint_D y\, \rho \dA . \]
  \pause
  \begin{example}[The region under $y = x^2$, $0 \le x \le 1$, with $\rho = 1$]
    $M = \int_0^1 x^2\, dx = \frac13$, \quad
    $\bar x = 3 \int_0^1 \int_0^{x^2} x\, dy\, dx = 3 \int_0^1 x^3\, dx = \frac34$, \\[0.2em]
    $\bar y = 3 \int_0^1 \int_0^{x^2} y\, dy\, dx = 3 \int_0^1 \frac{x^4}{2}\, dx = \frac{3}{10}$.
  \end{example}
  \pause
  {\small With $\rho = 1$ the centre of mass is the \textbf{centroid}: the average position of a point of $D$.\par}
\end{frame}

\begin{frame}{Probability}
  Random variables $X$, $Y$ have a \textbf{joint density} $p(x, y) \ge 0$, with
  $\iint_{\R^2} p \dA = 1$, if for regions $D$
  \[ P\bigl((X, Y) \in D\bigr) = \iint_D p(x, y) \dA . \]
  \pause
  \begin{example}[{$X$, $Y$ independent and uniform on $[0, 1]$}]
    Then $p = 1$ on the unit square and $0$ outside, so probabilities are areas:
    \begin{itemize}
      \item $P(X + Y \le 1) = \frac12$ (a triangle);
      \item $P(Y \le X^2) = \int_0^1 \int_0^{x^2} dy\, dx = \frac13$.
    \end{itemize}
  \end{example}
  \pause
  {\small Joint densities, expectations $E[g(X, Y)] = \iint g\, p \dA$, and independence are developed
  in Probability Theory and Statistics.\par}
\end{frame}

% =============================================================
\section{Summary and practice}

\begin{frame}{Summary}
  \begin{itemize}
    \item $\iint_R f \dA$: the limit of Riemann sums $\sum f(x_{ij}^*, y_{ij}^*)\,\Delta A$; for $f \ge 0$,
      the volume under the graph. Continuous functions are integrable.
    \item \textbf{Fubini}: for continuous $f$, the double integral is an iterated integral, in
      either order. Without continuity the orders can disagree.
    \item General regions: extend by zero. Type I: $\int_a^b \int_{g_1(x)}^{g_2(x)} \dots\, dy\, dx$;
      type II: $\int_c^d \int_{h_1(y)}^{h_2(y)} \dots\, dx\, dy$.
    \item Changing the order: sketch the region and describe it the other way. It can turn an
      impossible inner integral into an easy one.
    \item Area, average value, mass, centre of mass, probabilities.
  \end{itemize}
\end{frame}

\begin{frame}{Practice}
  \small
  \begin{enumerate}
    \item Estimate $\iint_R xy \dA$, $R = [0, 2]^2$, by a midpoint sum with $4$ squares, and compare
      with the exact value.
    \item Compute $\int_0^1 \int_0^2 (x + y)^2\, dy\, dx$.
    \item Compute $\iint_T (x + y) \dA$, where $T$ is the triangle with vertices $(0, 0)$, $(1, 0)$,
      $(0, 1)$.
    \item Find the area between $y = x^2$ and $y = 2x$.
    \item Compute $\displaystyle \int_0^\pi \int_x^\pi \frac{\sin y}{y}\, dy\, dx$ by changing the order.
    \item Compute $\displaystyle \int_0^1 \int_{\sqrt x}^1 \sin(y^3)\, dy\, dx$.
    \item $X$, $Y$ are independent and uniform on $[0, 1]$. Find $P(XY \le \frac12)$.
  \end{enumerate}
\end{frame}

\begin{frame}{Answers}
  \small
  \begin{enumerate}
    \item $\frac14 + \frac34 + \frac34 + \frac94 = 4$; \ exact: $\int_0^2 x\,dx \cdot \int_0^2 y\,dy = 4$.
      (The midpoint rule is exact for $xy$.)
    \item $\int_0^1 \bigl(2x^2 + 4x + \frac83\bigr)\, dx = \frac{16}{3}$.
    \item $\int_0^1 \int_0^{1-x} (x + y)\, dy\, dx = \frac16 + \frac16 = \frac13$.
    \item The curves meet at $x = 0$ and $x = 2$: \ $\int_0^2 (2x - x^2)\, dx = \frac43$.
    \item $\int_0^\pi \int_0^y \frac{\sin y}{y}\, dx\, dy = \int_0^\pi \sin y\, dy = 2$.
    \item The region is $0 \le y \le 1$, $0 \le x \le y^2$: \ $\int_0^1 y^2 \sin(y^3)\, dy = \frac{1 - \cos 1}{3}$.
    \item $1 - \iint_{xy > 1/2} 1 \dA = 1 - \int_{1/2}^1 \bigl(1 - \frac{1}{2x}\bigr) dx = \frac{1 + \ln 2}{2} \approx 0.847$.
  \end{enumerate}
\end{frame}

\begin{frame}{Next time}
  \textbf{Lecture 19 --- Triple integrals} (Unit 19)
  \begin{itemize}
    \item Integrals over boxes, and Fubini's theorem in three variables
    \item Integrals over general solids: choosing the order and the limits
    \item Mass, centre of mass, and probability in three dimensions
  \end{itemize}
  \vskip1em
  {\small
  Reading: Strang \& Herman, \emph{Calculus, Volume 3}, Sections 5.4 and 5.6.\\[0.4em]
  Questions: the module channel on Discord, or office hours,
  Monday--Friday, 10:00--17:00 (Baku time, UTC+4),
  \href{mailto:ali.khudiyev@wosm.academy}{ali.khudiyev@wosm.academy}.\par}
\end{frame}

\end{document}
